%% Question bank chapter file - generated by bank_tools.py build (do not edit while the bank is still being built:
%% it is regenerated from the per-test data in Question Bank\data). Every question is one bankq block:
%%   \begin{bankq}[id=..., source={...}, paper=.., qno=.., type=..., subtopic={...}, answer={...}, answer_src=key|solved,
%%                 working={...}, difficulty=1-10, diff_parts={K S I C F}, diff_note={...},
%%                 group=<passage id or empty>, also_in={duplicate copies in other tests}]
%% difficulty = min(10, K+S+I+C+F-1): K concepts 1-3, S steps 1-4, I insight 0-2, C computation 0-2, F format 0-1
%%   ...question LaTeX, options, figures (\bfig / \bfigi -> img/)...
%%   \end{bankq}
%% Passages shared by several questions are bankpassage blocks; their questions carry group=<passage id>.
\documentclass[10pt]{article}
\usepackage[a4paper,margin=1.8cm]{geometry}
\usepackage[T1]{fontenc}
\usepackage{amsmath,amssymb,mathtools}
\usepackage{graphicx}
\usepackage[export]{adjustbox}
\usepackage{xcolor,needspace,keyval}
\usepackage{enumitem}
\usepackage{booktabs,array,tabularx}
\usepackage[hidelinks]{hyperref}
\setlength{\parindent}{0pt}
\setlength{\parskip}{3pt}
\makeatletter
\define@key{bq}{id}{\def\bq@id{#1}}\define@key{bq}{source}{\def\bq@src{#1}}\define@key{bq}{paper}{\def\bq@paper{#1}}
\define@key{bq}{qno}{\def\bq@qno{#1}}\define@key{bq}{type}{\def\bq@type{#1}}\define@key{bq}{subtopic}{\def\bq@sub{#1}}
\define@key{bq}{answer}{\def\bq@ans{#1}}\define@key{bq}{answer_src}{\def\bq@asrc{#1}}\define@key{bq}{working}{\def\bq@work{#1}}
\define@key{bq}{group}{\def\bq@group{#1}}\define@key{bq}{also_in}{\def\bq@also{#1}}
\define@key{bq}{difficulty}{\def\bq@diff{#1}}\define@key{bq}{diff_parts}{\def\bq@dparts{#1}}\define@key{bq}{diff_note}{\def\bq@dnote{#1}}
\newenvironment{bankq}[1][]{\def\bq@work{}\def\bq@also{}\def\bq@group{}\def\bq@diff{?}\def\bq@dparts{}\def\bq@dnote{}\setkeys{bq}{#1}%
  \par\needspace{4\baselineskip}\medskip\noindent{\small\textbf{\bq@id}\hfill\color{gray}[\bq@src, Paper \bq@paper, Q\bq@qno]
  \bq@type\ $\cdot$ \bq@sub\ $\cdot$ \textbf{D\bq@diff}}\par\nopagebreak\relax}
  {\par{\small\color{teal}Answer: \textbf{\bq@ans}\ifx\bq@work\@empty\else\ (solved: \bq@work)\fi
  \ifx\bq@also\@empty\else\par\color{gray}Also in: \bq@also\fi}\par}
\define@key{bp}{id}{\def\bp@id{#1}}\define@key{bp}{source}{\def\bp@src{#1}}\define@key{bp}{questions}{\def\bp@qs{#1}}
\newenvironment{bankpassage}[1][]{\setkeys{bp}{#1}\par\needspace{4\baselineskip}\bigskip\noindent
  {\small\textbf{Passage \bp@id}\hfill\color{gray}[\bp@src; questions \bp@qs]}\par\nopagebreak\relax\itshape}{\par}
\makeatother
\newcommand{\bfig}[2]{\par\smallskip{\centering\includegraphics[width=#1pt,max width=0.62\linewidth,max height=0.32\textheight]{img/#2}\par}\smallskip}
\newcommand{\bfigi}[2]{\includegraphics[width=#1pt,max width=\linewidth,max height=0.32\textheight,valign=c]{img/#2}}
\newcommand{\bmissing}[1]{\textit{[figure missing in the source conversion: \detokenize{#1}]}}
\newenvironment{qoptions}{\begin{itemize}[leftmargin=1.8em,itemsep=1pt,topsep=2pt,label={}]}{\end{itemize}}
\newcommand{\opt}[2]{\item[(#1)] #2}
\providecommand{\fwstart}{}\providecommand{\fwend}{}\providecommand{\qfullsplit}{}
\newcolumntype{L}[1]{>{\raggedright\arraybackslash}p{#1}}
\newcolumntype{C}[1]{>{\centering\arraybackslash}p{#1}}
\newcolumntype{R}[1]{>{\raggedleft\arraybackslash}p{#1}}
\providecommand{\sectionhdr}[1]{}\providecommand{\subsectionhdr}[1]{}


\begin{document}
{\Large\bfseries Current Electricity --- Question Bank}\par\smallskip
{\small Physics chapter P10. Questions from all converted test papers, grouped by source series. Answers: printed key where it matches the paper, otherwise solved (answer\_src=solved, with working). Exact repeats across papers appear once, with also\_in listing the other copies.}\par\medskip\hrule

\section*{Score-I papers}
\begin{bankq}[id=SI12-P1-PHY-S1-Q01, source={Score-I - Test 12 (09 Mar 2022)}, paper=1, qno=1, type=multi, subtopic={Meters}, answer={B,C,D}, answer_src=key, difficulty=4, diff_parts={K1 S2 I1 C0 F1}, diff_note={concepts: series-parallel circuit with an ideal ammeter | steps: ammeter shorts part of the rheostat; current as a function of jockey position; minimum and maximum current; test each option}]
The circuit diagram shows a source of emf 10 V having an internal resistance of 2.5\,$\Omega$ connected to an ideal ammeter. The jockey can be slided to different points on the Rheostat whose total resistance is 10\,$\Omega$. What can be the reading of the ammeter?
\bfig{129.4}{SI12_p1_phy_s1_q01_1.png}
\begin{qoptions}
\opt{A}{1 A}
\opt{B}{2 A}
\opt{C}{3 A}
\opt{D}{4 A}
\end{qoptions}
\end{bankq}
\begin{bankq}[id=SI12-P1-PHY-S3-Q09, source={Score-I - Test 12 (09 Mar 2022)}, paper=1, qno=9, type=numerical, subtopic={Kirchhoff's laws}, answer={2}, answer_src=key, difficulty=2, diff_parts={K1 S1 I0 C0 F1}, diff_note={concepts: Kirchhoff's laws | steps: two loop equations; solve i1 and i2 and add the magnitudes}]
In the circuit indicated, find the current through both the batteries $i_1$ \& $i_2$. If $|i_1|+|i_2|=x$ then find the value of $\dfrac{5x}{6}$.
\bfig{129.3}{SI12_p1_phy_s3_q09_1.png}
\end{bankq}
\begin{bankpassage}[id=SI12-P2-PHY-S1-G5-6, source={Score-I - Test 12 (09 Mar 2022)}, questions={5-6}]
In the given circuits, all the ammeters are ideal but voltmeters may not be ideal. In the below column-1 $A_1,A_2,A_3$ are used to represent the magnitude of current through the ammeters $A_1,A_2,A_3$. In Column-2 $V_1$ and $V_2$ are readings of voltmeters. $R_{eq}$ is equivalent resistance for different situations.
\bfig{112.5}{SI12_p2_phy_s1_para5_6_1.png}
\begin{center}\footnotesize
\begin{tabular}{|p{0.6cm}|p{1.6cm}|p{0.5cm}|p{1.6cm}|p{0.5cm}|p{1.6cm}|}
\hline
 & \centering Column-1\newline $|A_2-A_1|$ & & \centering Column-2\newline $|V_2-V_1|$ & & \centering Column-3\newline $R_{eq}$\tabularnewline
\hline
(I) & $\dfrac{25}3$ & (i) & 30 & (P) & 5\\
(II) & 4 & (ii) & 3 & (Q) & 9\\
(III) & 3 & (iii) & 0 & (R) & 6\\
(IV) & 8 & (iv) & 4 & (S) & 4\\
\hline
\end{tabular}
\end{center}
\end{bankpassage}
\begin{bankq}[id=SI12-P2-PHY-S1-Q05, source={Score-I - Test 12 (09 Mar 2022)}, paper=2, qno=5, type=single, subtopic={Meters}, answer={C}, answer_src=key, difficulty=2, diff_parts={K1 S2 I0 C0 F0}, diff_note={concepts: circuits with non-ideal voltmeters | steps: equivalent resistance with voltmeter resistances; current; voltmeter readings; match the options}, group=SI12-P2-PHY-S1-G5-6]
If $R_1=3\,\Omega$, $R_2=6\,\Omega$. Resistance of voltmeters $V_1$ and $V_2$ are $6\,\Omega$ and $3\,\Omega$ respectively and $E=24$ V and both the switches are closed.
\begin{qoptions}
\opt{A}{(II) (ii) (R)}
\opt{B}{(IV) (iii) (R)}
\opt{C}{(II) (iii) (S)}
\opt{D}{(I) (iv) (Q)}
\end{qoptions}
\end{bankq}
\begin{bankq}[id=SI12-P2-PHY-S1-Q06, source={Score-I - Test 12 (09 Mar 2022)}, paper=2, qno=6, type=single, subtopic={Meters}, answer={C}, answer_src=key, difficulty=2, diff_parts={K1 S2 I0 C0 F0}, diff_note={concepts: circuits with non-ideal voltmeters | steps: equivalent resistance with voltmeter resistances; current; voltmeter readings; match the options}, group=SI12-P2-PHY-S1-G5-6]
$R_1=8\,\Omega$, $R_2=2\,\Omega$, $E=50$ V. Resistance of $V_1$ and $V_2$ is $12\,\Omega$ and $3\,\Omega$ respectively and both the switches are closed :-
\begin{qoptions}
\opt{A}{(I) (iv) (R)}
\opt{B}{(IV) (i) (Q)}
\opt{C}{(I) (i) (R)}
\opt{D}{(III) (iii) (S)}
\end{qoptions}

\noindent
\end{bankq}

\section*{TOAS papers}
\begin{bankq}[id=TOAS08-P1-PHY-S1-Q04, source={TOAS - Test 08 (15 May 2022)}, paper=1, qno=4, type=multi, subtopic={Kirchhoff's laws}, answer={A,B}, answer_src=key, difficulty=3, diff_parts={K1 S2 I0 C0 F1}, diff_note={concepts: Kirchhoff's laws in a multi-battery circuit | steps: loop equations; current in branch CD; power in the 10 V battery; heat in CH and the 20 V battery}]
In given circuit all batteries are ideal than:
\bfig{155.5}{TOAS08_p1_phy_s1_q04_1.png}
\begin{qoptions}
\opt{A}{Current in branch CD is 1.4 A from D to C.}
\opt{B}{10V battery is charging, consuming 6W power.}
\opt{C}{Rate of heat production in branch CH is 8.4 J/s.}
\opt{D}{In this circuit only 20V battery is discharging supplying power of 98W.}
\end{qoptions}
\end{bankq}
\begin{bankq}[id=TOAS08-P1-PHY-S3-Q01, source={TOAS - Test 08 (15 May 2022)}, paper=1, qno=1, type=numerical, subtopic={Meters}, answer={15.00}, answer_src=key, difficulty=2, diff_parts={K1 S1 I0 C0 F1}, diff_note={concepts: ammeter and voltmeter resistances | steps: voltage and current give the combined resistance; difference}]
In given circuit reading of ammeter and voltameter are $\dfrac23$A and $\dfrac{20}{3}$V respectively. If resistance of voltameter is $R_V$ and resistance of ammeter is $R_A$ then find $R_V-R_A$ (in $\Omega$).
\bfig{106.6}{TOAS08_p1_phy_s3_q01_1.png}
\end{bankq}
\begin{bankq}[id=TOAS08-P1-PHY-S3-Q05, source={TOAS - Test 08 (15 May 2022)}, paper=1, qno=5, type=numerical, subtopic={Ohm's law}, answer={70.00}, answer_src=key, difficulty=4, diff_parts={K1 S1 I2 C0 F1}, diff_note={concepts: infinite triangular resistor grid | steps: use symmetry of currents in an infinite lattice; equivalent resistance of adjacent nodes R/3 | I: symmetry argument for the infinite grid}]
In an infinite grid of triangular cell, resistance of each branch is $R=210\Omega$. The equivalent resistance of circuit between the terminal $A$ and $B$ is $x\Omega$. Find $x$.
\bfig{154.2}{TOAS08_p1_phy_s3_q05_1.png}
\end{bankq}
\begin{bankq}[id=TOAS08-P2-PHY-S1-Q02, source={TOAS - Test 08 (15 May 2022)}, paper=2, qno=2, type=multi, subtopic={Kirchhoff's laws}, answer={A,C,D}, answer_src=key, difficulty=5, diff_parts={K1 S3 I0 C1 F1}, diff_note={concepts: Kirchhoff's laws with several equal batteries | steps: assign loop currents; write the loop equations; solve; potential differences; power in a battery}]
Consider electric circuit shown in figure. All batteries are of emf 2V and all resistors are of $1\Omega$. Choose correct options :-
\bfig{158.0}{TOAS08_p2_phy_s1_q02_1.png}
\begin{qoptions}
\opt{A}{Current in resistor between $C$ \& $G$ is $\dfrac87$A.}
\opt{B}{$V_A-V_E$ is equal to $\dfrac{49}{7}$V.}
\opt{C}{$V_B-V_F$ is equal to $\dfrac{20}{7}$V.}
\opt{D}{Battery between $C$ \& $D$ supplies power of $\dfrac47$W.}
\end{qoptions}
\end{bankq}
\begin{bankq}[id=TOAS09-P1-PHY-S1-Q02, source={TOAS - Test 09 (30 May 2022)}, paper=1, qno=2, type=single, subtopic={Ohm's law}, answer={B}, answer_src=key, difficulty=1, diff_parts={K1 S1 I0 C0 F0}, diff_note={concepts: power transfer to a load | steps: power is maximum when R equals the internal resistance; read the peak and use P_max = E^2/4r}]
A battery of emf 48 V unknown internal resistance is connected to a rheostat whose resistance can be varied in a long range. A graph is plotted between power developed across rheostat and resistance of rheostat which is shown in figure. The internal resistance of the battery in $\Omega$ is :-
\bfig{118.2}{TOAS09_p1_phy_s1_q02_1.png}
\begin{qoptions}
\opt{A}{1.44}
\opt{B}{2.88}
\opt{C}{5.76}
\opt{D}{11.52}
\end{qoptions}
\end{bankq}
\begin{bankq}[id=TOAS09-P1-PHY-S1-Q04, source={TOAS - Test 09 (30 May 2022)}, paper=1, qno=4, type=single, subtopic={Potentiometer}, answer={D}, answer_src=key, difficulty=2, diff_parts={K1 S2 I0 C0 F0}, diff_note={concepts: potentiometer null condition | steps: potential gradient of the wire; null condition; solve R}]
A potentiometer wire AB is 100 cm long and has total resistance of 10$\Omega$. Find the value of unknown resistance R so that null point is obtained at a distance 40 cm from A.
\bfig{95.1}{TOAS09_p1_phy_s1_q04_1.png}
\begin{qoptions}
\opt{A}{1$\Omega$}
\opt{B}{2$\Omega$}
\opt{C}{3$\Omega$}
\opt{D}{4$\Omega$}
\end{qoptions}
\end{bankq}
\begin{bankq}[id=TOAS09-P1-PHY-S2-Q06, source={TOAS - Test 09 (30 May 2022)}, paper=1, qno=6, type=multi, subtopic={Meters}, answer={A,B,C,D}, answer_src=key, difficulty=3, diff_parts={K1 S2 I0 C0 F1}, diff_note={concepts: network of resistors tapped on a wire | steps: equivalent resistance from 400 V and 6 A; unknown R; current in the 25 ohm resistor; drop across 120 ohm}]
The resistance of wire PQ is 410 $\Omega$, which is tapped with the external wire as shown in the figure. The value of R is not known. If the potential difference between AB is 400 volts and the net current through the source battery is 6 ampere, then which of the following is/are correct :-
\bfig{149.2}{TOAS09_p1_phy_s2_q06_1.png}
\begin{qoptions}
\opt{A}{The resistance of unknown resistor R is 200 $\Omega$.}
\opt{B}{The equivalent resistance across AB is $\dfrac{200}{3}\,\Omega$.}
\opt{C}{The current in 25$\Omega$ resistor is $\dfrac{2}{5}$ ampere.}
\opt{D}{The potential drop across 120$\Omega$ resistor is 60 volts.}
\end{qoptions}
\end{bankq}
\begin{bankpassage}[id=TOAS09-P1-PHY-S3-G5-6, source={TOAS - Test 09 (30 May 2022)}, questions={5-6}]
In the circuit shown in figure, $R_1=R_2=R_3=3\,\Omega$ and e.m.f. of each cell is $E=4$V and negligible internal resistance. All ammeters are ideal.
\bfig{127.3}{TOAS09_p1_phy_s3_para5_6_1.png}
\end{bankpassage}
\begin{bankq}[id=TOAS09-P1-PHY-S3-Q05, source={TOAS - Test 09 (30 May 2022)}, paper=1, qno=5, type=numerical, subtopic={Kirchhoff's laws}, answer={1.33}, answer_src=key, difficulty=3, diff_parts={K1 S2 I0 C0 F1}, diff_note={concepts: circuit analysis from the passage | steps: equivalent resistance; branch currents; ammeter reading}, group=TOAS09-P1-PHY-S3-G5-6]
Reading of ammeter $A_1$ in ampere is:
\end{bankq}
\begin{bankq}[id=TOAS09-P1-PHY-S3-Q06, source={TOAS - Test 09 (30 May 2022)}, paper=1, qno=6, type=numerical, subtopic={Kirchhoff's laws}, answer={4.00}, answer_src=key, difficulty=3, diff_parts={K1 S2 I0 C0 F1}, diff_note={concepts: circuit analysis from the passage | steps: loop equations; potential at a and b; difference}, group=TOAS09-P1-PHY-S3-G5-6]
Potential difference (in Volt) between points a and b is
\end{bankq}
\begin{bankq}[id=TOAS09-P2-PHY-S1-Q01, source={TOAS - Test 09 (30 May 2022)}, paper=2, qno=1, type=single, subtopic={Kirchhoff's laws}, answer={B}, answer_src=key, difficulty=1, diff_parts={K1 S1 I0 C0 F0}, diff_note={concepts: Kirchhoff's laws | steps: node and loop equations; solve R}]
In the given circuit, find the value of resistance `R' :
\bfig{151.8}{TOAS09_p2_phy_s1_q01_1.png}
\begin{qoptions}
\opt{A}{20$\Omega$}
\opt{B}{25$\Omega$}
\opt{C}{15$\Omega$}
\opt{D}{10$\Omega$}
\end{qoptions}
\end{bankq}
\begin{bankq}[id=TOAS09-P2-PHY-S1-Q02, source={TOAS - Test 09 (30 May 2022)}, paper=2, qno=2, type=single, subtopic={Wheatstone bridge}, answer={B}, answer_src=key, difficulty=2, diff_parts={K1 S1 I1 C0 F0}, diff_note={concepts: balanced Wheatstone bridge | steps: identify the balanced bridge; combine the series-parallel branches | I: balanced bridge}]
In the given figure, $R=3\Omega$. Determine equivalent resistance between points A and B.
\bfig{142.1}{TOAS09_p2_phy_s1_q02_1.png}
\begin{qoptions}
\opt{A}{3 $\Omega$}
\opt{B}{2 $\Omega$}
\opt{C}{1.5 $\Omega$}
\opt{D}{6 $\Omega$}
\end{qoptions}
\end{bankq}
\begin{bankq}[id=TOAS09-P2-PHY-S1-Q03, source={TOAS - Test 09 (30 May 2022)}, paper=2, qno=3, type=single, subtopic={Ohm's law}, answer={A}, answer_src=key, difficulty=2, diff_parts={K1 S2 I0 C0 F0}, diff_note={concepts: series and parallel combinations of wire pieces | steps: resistance of each piece R/2n; series and parallel bundles; connect in series for the maximum}]
A piece of conducting wire of resistance $R$ is cut into $2n$ equal parts. Half the parts are connected in series to form a bundle and remaining half in parallel to form another bundle. These bundles are then connected to give the maximum resistance. The resistance of the combination is
\begin{qoptions}
\opt{A}{$\dfrac R2\left(1+\dfrac1{n^2}\right)$}
\opt{B}{$\dfrac R2\left(1+n^2\right)$}
\opt{C}{$\dfrac{R}{2(1+n^2)}$}
\opt{D}{$R\left(n+\dfrac1n\right)$}
\end{qoptions}
\end{bankq}
\begin{bankq}[id=TOAS09-P2-PHY-S1-Q04, source={TOAS - Test 09 (30 May 2022)}, paper=2, qno=4, type=single, subtopic={Kirchhoff's laws}, answer={B}, answer_src=key, difficulty=2, diff_parts={K1 S2 I0 C0 F0}, diff_note={concepts: current through a connecting wire | steps: equal potentials at both ends; split currents in each branch; current through AB}]
In figure, points A and B are connected by a perfectly conducting wire. Calculate the current through AB
\bfig{129.8}{TOAS09_p2_phy_s1_q04_1.png}
\begin{qoptions}
\opt{A}{2A}
\opt{B}{1 A}
\opt{C}{1.5 A}
\opt{D}{2.5 A}
\end{qoptions}
\end{bankq}
\begin{bankq}[id=TOAS09-P2-PHY-S1-Q05, source={TOAS - Test 09 (30 May 2022)}, paper=2, qno=5, type=single, subtopic={Potentiometer}, answer={D}, answer_src=key, difficulty=1, diff_parts={K1 S1 I0 C0 F0}, diff_note={concepts: potentiometer | steps: potential gradient from the driver cell; E2 from the balance length}]
AB is a potentiometer wire of length 100 cm. When a cell $E_1$ is connected across AC, where AC = 75 cm, no current flows from $E_2$. The internal resistance of the cell $E_1$ is negligible.
\bfig{151.8}{TOAS09_p2_phy_s1_q05_1.png}
Potential gradient along AB and value of $E_2$ respectively is.
\begin{qoptions}
\opt{A}{0.01 V cm$^{-1}$, 1.2 volt}
\opt{B}{0.03 Vcm$^{-1}$, 1.3 volt}
\opt{C}{0.04 Vm$^{-1}$, 1.0 volt}
\opt{D}{0.02 Vcm$^{-1}$, 1.5 volt}
\end{qoptions}
\end{bankq}
\begin{bankq}[id=TOAS09-P2-PHY-S1-Q06, source={TOAS - Test 09 (30 May 2022)}, paper=2, qno=6, type=single, subtopic={Kirchhoff's laws}, answer={C}, answer_src=key, difficulty=2, diff_parts={K1 S2 I0 C0 F0}, diff_note={concepts: ideal ammeter in a circuit | steps: equivalent circuit with the ammeter as a wire; branch currents; reading}]
Find the reading of the ideal ammeter connected in the given circuit. Assume that the cells have negligible internal resistance.
\bfig{152.0}{TOAS09_p2_phy_s1_q06_1.png}
\begin{qoptions}
\opt{A}{2.5 A}
\opt{B}{5 A}
\opt{C}{2 A}
\opt{D}{none of these}
\end{qoptions}
\end{bankq}
\begin{bankq}[id=TOAS09-P2-PHY-S2-Q09, source={TOAS - Test 09 (30 May 2022)}, paper=2, qno=9, type=numerical, subtopic={Ohm's law}, answer={5}, answer_src=key, difficulty=2, diff_parts={K1 S1 I0 C0 F1}, diff_note={concepts: steady current in a wire of varying cross-section | steps: current is constant so E is inversely proportional to area}]
Figure shows segment of a current carrying wire made of copper in steady state. If ratio of magnitudes of electric field at cross-section $S_1$ and $S_2$ is $5n$. Find value of n.
\bfig{153.7}{TOAS09_p2_phy_s2_q09_1.png}
\end{bankq}
\end{document}
